\documentclass[11pt,a4paper]{article}

\usepackage[utf8]{inputenc}
\usepackage[T1]{fontenc}
\usepackage{lmodern}
\usepackage{geometry}
\geometry{margin=2.5cm}
\usepackage{amsmath,amssymb,amsthm}
\usepackage{graphicx}
\usepackage{hyperref}
\usepackage{csquotes}

\title{%
Projective Dynamic Logo as a Foundational Research Programme:\\
Current Status, Open Problems, and Interface with the Theory of Objectivity
}

\author{C\'edric Laubscher\\
\small Independent researcher, Switzerland}

\date{February 2026}

\begin{document}

\maketitle

\begin{abstract}
The Projective Dynamic Logo (PDL) is a foundational research programme that conceptualises physical reality as a network of minimally closed logical structures defined on finite signed graphs.\cite{Laubscher2026IntroPDL,Laubscher2026FullPDL} Starting from four relational axioms and a coherence–optimisation principle, PDL yields a minimal $(4,6)$ stationary structure, interpreted as a prototypical electron, a hierarchical proton architecture with tightly constrained integer invariants, and a unified reinterpretation of several constants—such as $\hbar$, $\alpha$, $k_B$, and an effective gravitational coupling—in terms of coherence ratios.\cite{Laubscher2026FullPDL,Laubscher2026Alpha,Laubscher2026TopologicalReformulation} This position paper does not present new technical results. Its objective is to synthesise the current status of the PDL framework, to clearly differentiate between rigorously established structural results, tightly constrained conjectures, and calibrated elements, and to formulate a concise set of open problems that delineate the next stages of the programme. Particular emphasis is placed on the implications of the dialogue with the Theory of Objectivity (TO), which has sharpened the modal and ontological issues underlying PDL, especially with respect to boundaries, multi‑observer stabilisation, and the interpretation of coherence leakage as a transcendent substrate.\cite{Laubscher2026DialogueTO}
\end{abstract}

\section{Introduction and Aim}
The Projective Dynamic Logo (PDL) is a minimalist programme aimed at reconstructing physical reality from purely relational axioms formulated on finite signed graphs, without assuming a pre‑existing spacetime background, continuous fields, or primitive particle entities.\cite{Laubscher2026IntroPDL,Laubscher2026FullPDL} The central hypothesis is that existence and persistence can be reformulated as properties of closed, periodically pulsating relational patterns, governed by a coherence‑optimisation principle rather than by externally imposed dynamical laws.\cite{Laubscher2026FullPDL} Within this framework, a minimal $(4,6)$ structure emerges as the first admissible stationary closure, which is interpreted as a logical prototype of an electron at rest, while a more elaborate hierarchical configuration is employed to model the proton and to reinterpret several physical constants as internal coherence ratios.\cite{Laubscher2026IntroPDL,Laubscher2026Alpha}

The present article does not undertake a technical rederivation of the underlying constructions. Rather, it is intended as a concise position paper with three primary aims: (i) to provide a compact overview of the central structures of PDL and their current quantitative scope, (ii) to distinguish those components of the framework that can already be regarded as structurally robust from those that remain conjectural or empirically calibrated, and (iii) to present a systematic list of open problems which, if resolved, would either substantially corroborate or decisively undermine the programme.\cite{Laubscher2026FullPDL,Laubscher2026TopologicalReformulation} An additional objective is to elucidate how the ongoing dialogue with the Theory of Objectivity (TO) has sharpened the modal interpretation of PDL, particularly with respect to the roles of nothingness, boundaries, multi‑observer stabilisation, and a transcendent substrate.\cite{Laubscher2026DialogueTO}

The paper is structured as follows. Section~2 briefly reviews the core architecture of the PDL framework, including its relational axioms, the minimal $(4,6)$ block, the proton architecture, and the general strategy of interpreting constants as coherence ratios.\cite{Laubscher2026IntroPDL,Laubscher2026FullPDL} Section~3 summarises the conceptual implications of the TO–PDL dialogue and sketches informal points of compatibility with TO’s Seven Absolute Truths.\cite{Laubscher2026DialogueTO} Section~4 compiles the principal structural and quantitative results obtained thus far. Section~5 is devoted to explicit open problems and prospective research directions, spanning issues of combinatorial uniqueness as well as dynamical and phenomenological extensions. A brief conclusion argues that, taken together, these elements characterise PDL as a developing foundational research programme rather than a completed theory.

\section{Core Structure of the PDL Framework}

\subsection{Relational axioms and minimal closure}
In PDL, a \emph{logical configuration} is formalised as a finite signed graph $G = (V,E)$, where vertices represent informational entities and edges are assigned signs $s_{ij} \in \{+1,-1\}$ that encode binary relations of agreement or disagreement.\cite{Laubscher2026IntroPDL} On this discrete substrate, four relational axioms are postulated: minimal binary pulsation (existence instantiated as a repeatable two‑state alternation), triangular coherence (closed cycles of length three function as elementary closure motifs, with an associated penalty for each violated triangle), minimal completeness (all relational references required for persistence are constrained to remain internal to the structure), and logical optimisation (among all admissible configurations, the preferred ones are those that maximise an optimisation functional incorporating low violation, sufficient connectivity, and structural minimality).\cite{Laubscher2026IntroPDL,Laubscher2026FullPDL}

Subject to these constraints, it can be demonstrated that the first nontrivial stationary structure is realised by the complete signed graph on four vertices and six edges, denoted $(4,6)$. For this graph, there exists a sign assignment with zero fraction of triangle violations, together with a global sign inversion that implements a logical $2$‑cycle.\cite{Laubscher2026IntroPDL} This $(4,6)$ block is interpreted as the minimal self‑maintaining regime of existence and is proposed as a logical prototype for an electron at rest, with the reduced Planck constant reinterpreted as the minimal action associated with one internal coherence cycle.\cite{Laubscher2026FullPDL}

\subsection{Proton architecture and discrete invariants}
Beyond the elementary $(4,6)$ closure, PDL describes the proton as a hierarchically organised composite architecture rather than as a structureless effective degree of freedom.\cite{Laubscher2026FullPDL,Laubscher2026CombinatorialProton} In this framework, the protonal configuration comprises three quasi‑stationary valence cores, a dense relational sea, and a finite active surface mediating the coupling between the proton and external $(4,6)$ blocks.\cite{Laubscher2026FullPDL}

A combinatorial analysis governed by the same optimisation functional yields a set of discrete invariants: valence occupancies $n_u = 24$ and $n_d = 28$, a total valence relational content $r_{\text{val}} = 930$, a dynamical sea with $R_{\text{sea}} = 10\,087$ relations, and a total internal relational budget $R_{\text{tot}} = 11\,017$.\cite{Laubscher2026CombinatorialProton} This decomposition implies an approximate partition of internal coherence into $9\%$ quasi‑stationary and $91\%$ dynamical contributions, which is qualitatively consistent with QCD‑motivated pictures in which the dominant fraction of the proton mass originates from gluonic and sea components.\cite{Laubscher2026FullPDL} Within PDL, these integers are not introduced as independent parameters for separate phenomena; instead, they constitute a rigid combinatorial backbone that jointly constrains the proton’s internal closure, its coupling to electrons, and the formalisation of leakage mechanisms and surface‑defined constants.\cite{Laubscher2026Alpha,Laubscher2026TopologicalReformulation}

\subsection{Constants as coherence ratios}
Within this discrete relational framework, several dimensionless physical constants are reinterpreted as coherence ratios rather than treated as primitive empirical inputs.\cite{Laubscher2026FullPDL} The reduced Planck constant is construed as quantifying the minimal coherence expenditure associated with a complete pulsation of the $(4,6)$ regime.\cite{Laubscher2026IntroPDL} The fine‑structure constant $\alpha$ is modeled as a surface‑to‑volume coherence fraction, namely the proportion of the proton’s total coherence budget that is effectively projected onto an active surface available for coherent coupling to an external electron block.\cite{Laubscher2026Alpha} 

In a dedicated derivation, this conceptualization yields a closed‑form expression
\[
\alpha^{-1}_{\mathrm{PDL}} = \mu / (9 \varphi^{2} r_{\mathrm{val}}^{2}),
\]
where $\mu = m_p/m_e$ denotes the proton–electron mass ratio, $r_{\mathrm{val}} = 930$ represents the valence content, and $\varphi$ is the golden ratio. This expression reproduces the empirical value of $\alpha$ with a relative accuracy of order $10^{-4}$, without the introduction of additional continuous fitting parameters.\cite{Laubscher2026Alpha,Laubscher2026GoldenRatioSurface} 

Further notes generalize this approach to Boltzmann’s constant and to an effective gravitational coupling by introducing a small coherence‑leakage parameter $\varepsilon$ associated with composite nucleonic closures, and by interpreting gravitation as a modulation of coherence cost within an emergent relational metric.\cite{Laubscher2026FullPDL,Laubscher2026TopologicalReformulation,Laubscher2026Exponent18} Across these constructions, the objective is not precise numerical prediction per se, but rather to assess the extent to which a single discrete architecture, characterized by a limited set of structural parameters, can account for multiple constants as distinct manifestations of a unified underlying coherence structure.\cite{Laubscher2026FullPDL}

\section{Impact of the TO–PDL Dialogue}

\subsection{Conceptual clarifications triggered by TO}

The Theory of Objectivity (TO) addresses foundational questions from a modal–axiomatic standpoint, organised around seven Absolute Truths that any admissible universe must satisfy in order to sustain temporally persistent existence.\cite{Laubscher2026DialogueTO} The interaction between PDL and TO has therefore not been merely interpretive; rather, it has operated as an external modal probe into the ontological and logical commitments of PDL.\cite{Laubscher2026DialogueTO} Subjecting PDL to TO’s axiomatic constraints has enforced a sharper demarcation between, first, those claims in PDL that are to be regarded as logically necessary; second, those that are better construed as structured yet revisable conjectures; and third, those that are calibrated inputs grounded in empirically fixed constants.\cite{Laubscher2026IntroPDL,Laubscher2026FullPDL} This stratification has, in turn, clarified which components of the framework warrant the designation “foundational” and which remain at the level of specific model construction.

Several conceptual domains have been refined as a result of this exchange. First, TO’s insistence on taking nothingness as a primitive mathematical essence has motivated a more explicit reconstruction of the PDL axioms from conditions on reproducible distinction in a pre‑metric setting, as opposed to relying on an implicit background of graph‑theoretic formalism.\cite{Laubscher2026IntroPDL,Laubscher2026DialogueTO} Second, TO’s focus on boundaries as ontologically primary entities has led to a more rigorous treatment of active surfaces in PDL, with a clear differentiation between structural boundaries and purely representational cuts.\cite{Laubscher2026FullPDL,Laubscher2026Alpha,Laubscher2026DialogueTO} Third, TO’s multi‑observer axiom has motivated the introduction of a coexistence topology and a coherence‑cost pseudometric on the space of closures, thereby opening a pathway towards formalising existence in PDL as a matter of compatibility across multiple relational environments, rather than as a solely intra‑systemic property.\cite{Laubscher2026TopologicalReformulation,Laubscher2026DialogueTO} Finally, TO’s notion of a transcendent substance has sharpened the interpretive question of whether PDL’s coherence‑leakage parameter should be understood as encoding a genuinely exported substrate, or instead as a bookkeeping device for residual inconsistency.\cite{Laubscher2026FullPDL,Laubscher2026Exponent18,Laubscher2026TopologicalReformulation}

\subsection{Compatibilities with A1--A7 (informal overview)}
At an informal level, several systematic correspondences between PDL and TO’s Seven Absolute Truths can be discerned.\cite{Laubscher2026DialogueTO} Regarding A1 (nothingness), PDL’s point of departure—a pre‑metric logical field in which only reproducible distinctions can ground existence—may be construed as a formalisation of “nothingness” supplemented by minimal persistence conditions, without any prior assumption of spacetime structure or material substrate.\cite{Laubscher2026IntroPDL,Laubscher2026FullPDL} A2 (aura and uniqueness) aligns with the PDL procedure of individuating elements via relational invariants: the $(4,6)$ structure, the proton’s integer signature $(24,28,930,10\,087,11\,017)$, and the recurrence of these invariants across distinct constructions provide candidates for non‑conventional “auras” that are more than arbitrary identifiers.\cite{Laubscher2026FullPDL,Laubscher2026CombinatorialProton}

A4 (boundary necessity) admits a natural analogue in PDL’s active surfaces and closed logical regimes; distinctions in PDL are intrinsically associated with finite boundaries across which coherence is exchanged, and quantities such as $\alpha$ are explicitly defined in relation to these boundaries.\cite{Laubscher2026Alpha,Laubscher2026FullPDL} A5 (multi‑observer validation), while not yet axiomatic within PDL, is reflected in the requirement that closures be compatible with multiple surrounding closures and in the emerging coexistence topology. This suggests a graph‑theoretic reformulation of TO’s requirement for at least two independent “observers.”\cite{Laubscher2026TopologicalReformulation,Laubscher2026DialogueTO} A7 (transcendent substance) can be provisionally associated with the coherence‑leakage parameter $\varepsilon$ and its hierarchical filtration: a minimal, irreducible defect of closure that is not eliminable within any given quantum system, but instead appears as a global modulation of coherence cost and emergent metric curvature.\cite{Laubscher2026FullPDL,Laubscher2026Exponent18,Laubscher2026TopologicalReformulation} In contrast, A3 (infinity as non‑element) and A6 (compositional ancestry) currently find only partial counterparts in PDL and delineate domains in which the generative hierarchy and completion properties of closures would need to be articulated more explicitly.\cite{Laubscher2026DialogueTO}

\section{Current Achievements}

\subsection{Structural results on particles and architecture}
On the structural side, PDL has reached a level of development at which several core components are now specified with sufficient mathematical and physical precision to permit rigorous scrutiny.\cite{Laubscher2026FullPDL,Laubscher2026IntroPDL} The minimal $(4,6)$ structure is no longer introduced as a mere postulate but is instead characterised as the first configuration that satisfies the axioms of minimal pulsation, triangular coherence, minimal completeness, and optimisation: it is realised as a complete graph on four vertices, in which all triangular subgraphs are coherent with respect to an appropriate sign assignment, are covered by a connected network of coherent cycles, and admit a global sign‑flip pulsation.\cite{Laubscher2026IntroPDL} While a full uniqueness theorem remains outstanding, this specification already yields a concrete and well‑delimited object for combinatorial investigation. Similarly, the proton architecture has been elevated from an informal schematic representation to a rigorously defined hierarchical structure subject to integer constraints, closure relations, and a precisely specified stationary/dynamical decomposition.\cite{Laubscher2026FullPDL,Laubscher2026CombinatorialProton}

These constructions jointly instantiate a coherent relational account of “microscopic matter” within PDL: electrons are modelled as minimal self‑pulsating closures, protons as minimal hierarchical composites assembled from such closures together with a relational background sea, and more complex aggregates as networks of mutually compatible closures interconnected via active surfaces.\cite{Laubscher2026FullPDL} Crucially, an identical combinatorial substrate supports multiple physical roles: the integer data specifying the proton enter not only its internal stability conditions but also the surface degrees of freedom relevant for interaction couplings, the definition of coherence leakage, and the emergent metric structure.\cite{Laubscher2026Alpha,Laubscher2026TopologicalReformulation} This systematic reuse of internal structural information is a primary reason why PDL can be regarded as a unified theoretical framework rather than as a mere collection of isolated toy models.

\subsection{Quantitative correspondences for constants}
From a quantitative perspective, PDL yields several non‑trivial correspondences between its discrete combinatorial architecture and empirically measured fundamental constants. The most developed example concerns the fine‑structure constant $\alpha$, where a structural hypothesis identifies $\alpha^{-1}$ with a surface‑coherence ratio that involves the proton–electron mass ratio, the valence‑relation multiplicity, and a golden‑ratio–based optimisation factor. This leads to a theoretical estimate $\alpha^{-1}_{\mathrm{PDL}} \approx 137.02$, in agreement with the experimentally determined value at the level of a few parts in $10^{-4}$, obtained without introducing additional continuous fit parameters.\cite{Laubscher2026Alpha,Laubscher2026GoldenRatioSurface} Analogous, though currently less precise, constructions have been proposed for Boltzmann’s constant $k_B$ and for an effective gravitational coupling, in which a small coherence‑leakage parameter $\varepsilon$ together with a hierarchical exponent is employed to relate nucleonic closure defects to macroscopic metric curvature.\cite{Laubscher2026FullPDL,Laubscher2026Exponent18}

These numerical agreements do not amount to derivations in the strict logical or deductive sense; rather, they rely on structured, model‑level hypotheses, such as the identification of $\varepsilon$ with a minimal closure defect and the interpretation of certain ratios as effective coupling constants in familiar dynamical equations.\cite{Laubscher2026FullPDL,Laubscher2026TopologicalReformulation} Nevertheless, the observation that a single discrete architecture, characterized by a small set of structural parameters, can reproduce several independent orders of magnitude and dimensionless quantities is statistically non‑trivial. This indicates that the combinatorial choices encoded in PDL are not arbitrary and that the framework is sufficiently constrained to be empirically testable: tightening the axioms, exploring alternative integer assignments, or increasing experimental precision could, in principle, falsify the current proposals.\cite{Laubscher2026FullPDL,Laubscher2026Alpha}

\section{Open Problems and Future Directions}

\subsection{Foundational and combinatorial open problems}
A number of foundational issues remain unresolved and delineate several well‑defined research programmes.\cite{Laubscher2026FullPDL,Laubscher2026CombinatorialProton} 

First, the uniqueness of the $(4,6)$ block within the PDL axiomatic framework has not yet been established. A rigorous proof or disproof that no smaller, or otherwise distinct, combinatorial structure can satisfy the same closure and optimisation conditions would substantially clarify whether the $(4,6)$ regime is a strictly necessary structural primitive or merely one admissible choice among others.\cite{Laubscher2026IntroPDL}

Second, the combinatorial rigidity of the proton architecture—specified by the integer tuple $(24,28,930,10\,087,11\,017)$ together with the $9/91$ stationary/dynamical partition—currently rests on optimisation heuristics and compatibility with phenomenological data, rather than on an exhaustive combinatorial analysis.\cite{Laubscher2026CombinatorialProton} A systematic exploration of the surrounding space of admissible integer configurations is therefore required to determine whether the proposed architecture is indeed uniquely selected (up to isomorphism) by the PDL constraints, or whether it should instead be regarded as one element of a larger equivalence class of viable configurations.

Third, the definition and quantitative determination of the leakage parameter $\varepsilon$ call for a fully combinatorial characterisation. At present, $\varepsilon$ is introduced conceptually as a minimal, irreducible closure defect in hierarchical structures, but its numerical value remains partially calibrated against gravitational data.\cite{Laubscher2026FullPDL,Laubscher2026TopologicalReformulation} A satisfactory foundational account would derive $\varepsilon$ from intrinsic graph‑theoretic invariants of nucleonic architectures and only at a later stage compare the resulting prediction with experimental determinations of $G$.\cite{Laubscher2026TopologicalReformulation,Laubscher2026Exponent18}

A closely related open problem concerns the exponent $18$ appearing in the hierarchical amplification of leakage. Although a conceptual picture has been proposed in terms of eighteen coherence filters organised into three blocks of six, there is as yet no rigorous proof that $18$ is the minimal exponent compatible with all structural and phenomenological constraints.\cite{Laubscher2026Exponent18}

\subsection{Interface with TO: modal status and multi-observer axioms}
On the TO side, several outstanding problems concern the modal status of PDL structures.\cite{Laubscher2026DialogueTO} A central question is the formalisation of a PDL analogue of TO’s axiom A5 (multi‑observer stabilisation). The coexistence topology together with the coherence‑cost metric suggests the following heuristic principle: a closure should be accorded “full existence” only if it admits embeddings into at least two independent compatibility neighbourhoods. However, this requirement has not yet been elevated to the status of an explicit axiom or derived theorem.\cite{Laubscher2026TopologicalReformulation,Laubscher2026DialogueTO} Systematically investigating the implications of such a condition—for the taxonomy of closures, for the admissibility of specific architectural configurations, and for the operational interpretation of measurement—constitutes a substantial avenue for future research.

A further issue is whether the relational invariants of PDL can be legitimately promoted to TO‑style “auras.” To justify this upgrade, one would need to demonstrate that certain integer‑valued signatures and structural motifs are entailed by the axioms and optimisation principles themselves, rather than being artefacts of underdetermined modelling choices.\cite{Laubscher2026CombinatorialProton,Laubscher2026DialogueTO} In addition, the ontological status of both the leakage parameter and the emergent metric remains to be clarified. If $\varepsilon$ is to function as a transcendent substrate in TO’s sense, it must be shown to encode exported coherence that genuinely lies beyond any fixed quantum domain, rather than merely serving as a parametrisation of internal dissipative processes.\cite{Laubscher2026FullPDL,Laubscher2026TopologicalReformulation}

\subsection{Dynamical and phenomenological extensions}
In addition to the static architecture and constant–level correspondences, PDL encounters a number of unresolved issues in the domains of dynamics and phenomenology.\cite{Laubscher2026SchrodingerSketch,Laubscher2026SchrodingerPDL} An important objective is to derive, at least in an appropriate coarse‑grained regime, a linear wave equation analogous to the Schrödinger equation directly from the underlying graph dynamics, rather than postulating Schrödinger dynamics as an external effective law into which PDL‑defined constants are merely substituted.\cite{Laubscher2026SchrodingerSketch} Initial analyses indicate potential derivations based on discrete Laplacians and flux‑balance conditions on signed graphs; however, a fully explicit derivation, including a transparent specification of the approximations involved and the associated domains of validity, remains to be developed.\cite{Laubscher2026SchrodingerSketch,Laubscher2026TopologicalReformulation}

From a phenomenological perspective, generalizing the current constructions beyond the hydrogen atom and single‑proton systems would constitute a significant test. Implementing the active‑surface and leakage framework for other nuclei, for hydrogen‑like ions, or for simple multi‑electron systems could clarify whether the coherence‑ratio description exhibits a consistent scaling behavior with increasing structural and dynamical complexity.\cite{Laubscher2026Alpha,Laubscher2026CombinatorialProton} Moreover, achieving a more systematic unification of the treatments of $\alpha$, $k_B$, and the gravitational coupling within a single, internally coherent scheme—potentially anchored in a refined characterization of $\varepsilon$ and the role of the golden ratio—remains a central objective if PDL is to be advanced as a genuinely unified foundational framework.\cite{Laubscher2026FullPDL,Laubscher2026Alpha,Laubscher2026TopologicalReformulation}

\section{Conclusion}
The Projective Dynamic Logo (PDL) currently occupies an intermediate epistemic status between a speculative research proposal and a fully developed theoretical framework.\cite{Laubscher2026FullPDL} On the one hand, it exhibits a sharply delineated axiomatic core formulated in terms of signed graphs, a concrete minimal $(4,6)$ configuration, a detailed proton architecture characterized by integer invariants, and several quantitatively non‑trivial correspondences between discrete coherence ratios and empirically measured physical constants.\cite{Laubscher2026IntroPDL,Laubscher2026Alpha,Laubscher2026TopologicalReformulation} On the other hand, key theorems remain unproven, several parameters continue to depend on empirical calibration rather than derivation, and the dynamical as well as modal components of the framework are, at present, only partially articulated.\cite{Laubscher2026FullPDL,Laubscher2026Exponent18,Laubscher2026DialogueTO} The ongoing dialogue with the Theory of Objectivity has been pivotal in elucidating both the strengths and limitations of PDL, and it indicates that the latter is best regarded as an active foundational research programme with explicitly formulated conjectures and well‑defined decision points.\cite{Laubscher2026DialogueTO} It is hoped that this concise overview will enable prospective collaborators and critics to identify where their respective areas of expertise may most productively intersect with the open problems delineated above.

\section*{Acknowledgements}
The author wishes to acknowledge Marco Silva and the Theory of Objectivity community for their detailed and intellectually demanding feedback on earlier formulations of the PDL framework, feedback that has significantly contributed to refining the modal and structural objectives of the programme.\cite{Laubscher2026DialogueTO} Informal discussions with colleagues working in the foundations of physics, together with the constructive pressure exerted by external critiques, have likewise been instrumental in sharpening several of the questions presented here as open problems. The author further acknowledges the pragmatic assistance of advanced AI tools in drafting and organizing portions of recent manuscripts, while retaining sole responsibility for all scientific claims, conjectures, and any remaining errors.\cite{Laubscher2026FullPDL}

\bibliographystyle{plain}
\bibliography{PDL_position_paper}

\end{document}
